\SourceLeaf{270}%
folglich, wenn man mit Hülfe von 5 entweder \(\mathrm{d}U\) oder \(\mathrm{d}w\) eliminirt,
\[
\frac{\mathrm{d}\log m}{\mathrm{d}w}=\frac{(1-ee)(\alpha\cos U-\cos w)}{(1-ee\cos w^2)\sin w} \quad .\ .\ .\ .\ .\ .\ .\ .\ .\ .\ .\ .\quad(6)
\]
\[
\frac{\mathrm{d}\log m}{\mathrm{d}U}=\operatorname{cotg}U-\frac{\cos w}{\alpha\sin U}=\frac{\alpha\cos U-\cos w}{\alpha\sin U} \quad .\ .\ .\ .\ .\ .\ .\quad(7)
\]

Durch eine nochmalige Differentiation der Gleichung 7 erhält man
\[
\begin{aligned}
\frac{\mathrm{dd}\log m}{\mathrm{d}U^2}&=-\frac{1}{\sin U^2}+\frac{\cos U\cos w}{\alpha\sin U^2}+\frac{\sin w}{\alpha\sin U}\mathord{.}\frac{\mathrm{d}w}{\mathrm{d}U}\\
&=-\frac{1}{\sin U^2}+\frac{\cos U\cos w}{\alpha\sin U^2}+\frac{(1-ee\cos w^2)\sin w^2}{\alpha\alpha(1-ee)\sin U^2} \quad .\ .\ .\ .\ .\ .\quad(8)
\end{aligned}
\]

Soll nun für eine bestimmte Breite (Normalbreite) der Werth von \(m\) der Einheit gleich werden, für andere Breiten hingegen nur um Grössen der dritten Ordnung von 1 abweichen, die Breitenunterschiede als Grössen erster Ordnung betrachtet, so muss, wenn die Normalbreite auf dem Ellipsoid mit \(P\), die entsprechende auf der Kugel mit \(Q\) bezeichnet wird, für \(w=90^0-P\), \(U=90^0-Q\) in Gemässheit der Gleichungen 4. 7, 8 sein:
\[
A=\frac{a\cos P}{\alpha\cos Q\BareRoot{(1-ee\sin P^2)}} \quad .\ .\ .\ .\ .\ .\ .\ .\ .\ .\ .\ .\ .\ .\quad(9)
\]
\[
\alpha\sin Q=\sin P \quad .\ .\ .\ .\ .\ .\ .\ .\ .\ .\ .\ .\ .\ .\ .\ .\ .\ .\ .\quad(10)
\]
\[
0=1-\frac{\sin P\sin Q}{\alpha}-\frac{(1-ee\sin P^2)\cos P^2}{\alpha\alpha(1-ee)}
\]
oder, wenn man in letzterer Gleichung für \(\sin Q\) seinen Werth aus 10 substituirt,
\[
\alpha\alpha=1+\frac{ee\cos P^4}{1-ee} \quad .\ .\ .\ .\ .\ .\ .\ .\ .\ .\ .\ .\ .\ .\ .\ .\quad(11)
\]

Durch diese Gleichung ist demnach \(\alpha\) gegeben, sobald für \(P\) ein bestimmter Werth gewählt ist; \(Q\) kann sodann durch Gleichung 10, und \(A\) durch Gleichung 9 bestimmt werden; endlich ergibt sich \(k\) durch die Substitution von \(w=90^0-P\), \(U=90^0-Q\) in der allgemeinen Gleichung 3, nemlich
\[
k=\frac{\operatorname{tang}(45^0+\frac{1}{2}P)^\alpha}{\operatorname{tang}(45^0+\frac{1}{2}Q)}\mathord{.}(\tfrac{1-e\sin P}{1+e\sin P})^{\frac{1}{2}\alpha e} \quad .\ .\ .\ .\ .\ .\ .\ .\quad(12)
\]

\Article{4}
\indent Die Berechnung der Constanten \(A,\alpha,k\) und der Normalbreite auf der Kugel \(Q\) aus \(P\) und \(e\) wird man, da alle diese Grössen wie Grundlagen für die
